\section{The microscopic concentration charge and the exact-label endpoint}
\label{sec:v84-concentration}
The microscopic limit has an additional rigidity: shrinking the outer
width cannot gradually reduce a nonzero concentration charge. It is
useful to state this directly, rather than to regard the tail height
as an arbitrary modulus of that outer width.

Fix $0<\eta<1/2$, write $\kappa=\kappa_\eta$, and use exactly the
actual source $v^\eta$ and offset convention of
\eqref{eq:v83-dyadic-tail}. For a family of nonnegative sources $f_m$
with the same parameter and label indices, abbreviate
\[
 \mathcal H_M(f)=\limsup_m\sup_{R,n,k}m^2
                \esssup_{t\in\mathcal C_{m,R,n,k}(M)}f_m(t).
\]
This functional is monotone and subadditive. It is not assumed to
commute with infinite sums. Define the canonical concentration charge
by
\begin{equation}\label{eq:v84-canonical-charge}
 \mathfrak A_{\eta,M}=\lim_{\delta\downarrow0}
                     \mathcal H_M(\mathsf P_{\eta,\delta}v^\eta).
\end{equation}
The limit exists monotonically in $[0,\infty]$. The complete finite-count
potential is inside $\mathcal H_M$; no density is truncated after
taking its collision limit.

\begin{theorem}[Outer-scale rigidity of the microscopic tail]
\label{thm:v84-charge}
For every fixed sufficiently small $r>0$,
\begin{align}
 \mathfrak T_{\eta,M}(r)
       &\le\mathfrak A_{\eta,M}
                           \le2^\kappa\mathfrak T_{\eta,M}(r),
                                            \label{eq:v84-charge-comparison}\\
 \mathfrak T_{\eta,M}(2^{-j}r)
       &=\mathfrak T_{\eta,M}(r)\quad(j\in\mathbb Z),
                                            \label{eq:v84-dyadic-invariance}\\
 \mathcal H_M(\mathsf P_{\eta,r}v^\eta)
       &\le C_\eta r^{1-\kappa}
                                +2^\kappa\mathfrak A_{\eta,M}.
                                            \label{eq:v84-charge-bound}
\end{align}
In \eqref{eq:v84-dyadic-invariance} both widths are required to be
in the stated domain. The following conditions are equivalent:
$\mathfrak A_{\eta,M}=0$; the dyadic tail height vanishes at one
positive width; it vanishes at every positive width; and
$\mathfrak T_{\eta,M}(r)\to0$ as $r\downarrow0$.
\end{theorem}
\begin{proof}
Put $\delta_N=2^{-N}r$. At each finite collision count, the annuli
in $T_{\eta,r,N}$ partition $(0,\delta_N]$, and hence
\[
 T_{\eta,r,N}\le\mathsf P_{\eta,\delta_N}v^\eta
                                      \le2^\kappa T_{\eta,r,N}.
\]
Apply $\mathcal H_M$ before passing to $N\to\infty$. Since
$\delta_N\downarrow0$, monotonicity identifies the middle limit
with \eqref{eq:v84-canonical-charge} and proves
\eqref{eq:v84-charge-comparison}. The exact finite-count identity
$T_{\eta,r/2,N}=T_{\eta,r,N+1}$ proves dyadic invariance. Iteration
also proves the assertion for negative integers within the domain.

For a fixed $N$, split the dyadic sum at $N$. The head consists of
finitely many intervals fixed before the collision limit. The
fixed-width source input bounds its ordered height by
$C_\eta\sum_{j<N}(2^{-j}r)^{1-\kappa}\le C_\eta r^{1-\kappa}$.
Keep the entire remaining tail inside $\mathcal H_M$, and only then
send $N\to\infty$. Use \eqref{eq:v84-charge-comparison} to obtain
\eqref{eq:v84-charge-bound}. Finally, a zero tail at one width forces
the charge to be zero, and a zero charge forces all tails to vanish.
If the tails tend to zero with $r$, the lower bound
$\mathfrak T_{\eta,M}(r)\ge2^{-\kappa}\mathfrak A_{\eta,M}$
forces the same conclusion. No claim of equality at arbitrary
nondyadically related widths is needed.
\end{proof}

\begin{remark}[The remaining estimate is uniform integrability]
\label{rem:v84-charge-scope}
At any one fixed count, local integrability of $v^\eta$ implies
that its potential is finite for almost every roof and its truncated
potential tends to zero there. This fact alone says nothing about
\eqref{eq:v84-canonical-charge}. The charge retains the essential
supremum and collision limsup before removing the microscopic width.
The finite-count envelope of Proposition~\ref{prop:v83-envelope}
is one sufficient way to prove that charge zero; it is not an
additional theorem about the Lorentz source. In particular choosing
a smaller outer reconstruction width does not pay a nonzero charge.
\end{remark}

\begin{theorem}[The endpoint after residence and boundary transfer]
\label{thm:v84-endpoint}
For the unchanged scalar positive remainder, every fixed
$0<\eta<1/2$, all sufficiently large fixed $B$ and sufficiently
small previously fixed $r>0$,
\begin{align}
 \mathcal E_M\le{}& C_MB^{-1/192}
      +C\mathfrak J_M(r)+C\mathfrak K_M(r)
      +C_\eta\mathfrak A_{\eta,M}\notag\\
      &+C_\eta r^{1/32}
              +C_\Theta(1+\varepsilon(B)^{-1})\sqrt r.
                                            \label{eq:v84-endpoint}
\end{align}
Thus the original pointwise assertion follows if the canonical
microscopic charge is zero and the paired nongrazing and short
critical surface heights tend to zero. These conditions have not
been inferred from local mass or weak-$L^q$ bounds.
\end{theorem}
\begin{proof}
The source-level comparison in the proof of
Theorem~\ref{thm:v83-endpoint}, prior to separating its dyadic tail,
gives
\[
 \mathcal E_M\le C_MB^{-1/192}+C\mathfrak F_M^{\rm all}(r)
  +C_\eta\mathcal H_M(\mathsf P_{\eta,r}v^\eta)
  +C_\Theta(1+\varepsilon(B)^{-1})\sqrt r.
\]
Apply \eqref{eq:v84-surface-reduction} and
\eqref{eq:v84-charge-bound}. Since $1-\kappa>1/2>1/32$ and
$r<1$, the explicit potential contribution is bounded by
$C_\eta r^{1/32}$. This proves \eqref{eq:v84-endpoint}.

For the stated sufficient conclusion choose
$r(B)=c\varepsilon(B)^\gamma$ with fixed $c>0$ and $\gamma>2$.
At each $B$ the width, its positive cutoff $c_{r(B)}$, and all
auxiliary bands are fixed before the collision limit. Enlarge the
auxiliary bands afterwards and take $B\to\infty$ last. The three
proved explicit errors vanish, and the three stated conditions pay
the remaining nonnegative quantities. No arithmetic class has been
removed or normalized in this argument.
\end{proof}

The controlled surface height in Theorem~\ref{thm:v84-long-height}
and the controlled optical transfer in Theorem~\ref{thm:v84-paired-transfer}
are genuine estimates on the original physical source. The unresolved
geometric quantities in \eqref{eq:v84-endpoint} are now supported
on maximal regular segments of length at most $2r$, either with two
nongrazing true ends or with outgoing gradient of order $\sqrt r$.
The original pointwise theorem, rather than a local-norm substitute,
remains the governing assertion. Scalar closure must still precede
all same-roof path and likelihood normalizations.
